% Standalone problem document, generated from the adjacent README.md.
% Compile with XeLaTeX. Mathematical content is maintained in Markdown.
\documentclass[11pt,a4paper]{article}
\usepackage[fontsize=10.5pt]{scrextend}
\usepackage[margin=22mm,headheight=15pt,headsep=9mm,footskip=11mm]{geometry}
\usepackage{fontspec}
\setmainfont{texgyrepagella}[Extension=.otf,UprightFont=*-regular,BoldFont=*-bold,ItalicFont=*-italic,BoldItalicFont=*-bolditalic]
\setsansfont{texgyreheros}[Extension=.otf,UprightFont=*-regular,BoldFont=*-bold,ItalicFont=*-italic,BoldItalicFont=*-bolditalic]
\setmonofont{texgyrecursor}[Extension=.otf,UprightFont=*-regular,BoldFont=*-bold,ItalicFont=*-italic,BoldItalicFont=*-bolditalic,Scale=MatchLowercase]
\usepackage{amsmath,amssymb}
\usepackage{unicode-math}
\setmathfont{texgyrepagella-math.otf}
\usepackage{xcolor,microtype,enumitem,needspace,fancyhdr}
\usepackage{bookmark}
\definecolor{ink}{HTML}{17344B}
\definecolor{accent}{HTML}{197D80}
\definecolor{muted}{HTML}{596773}
\hypersetup{colorlinks=true,linkcolor=accent,urlcolor=accent,pdftitle={MI-24 - Schatten
norm complement involving Lin's positive definite
quantity},pdfauthor={Open Problems in Numerical Linear Algebra}}
\urlstyle{same}
\setlength{\parindent}{0pt}
\setlength{\parskip}{5pt plus 1pt minus 1pt}
\linespread{1.04}
\setlength{\emergencystretch}{3em}
\widowpenalty=10000
\clubpenalty=10000
\displaywidowpenalty=10000
\allowdisplaybreaks[1]
\setlist{nosep,leftmargin=1.5em}
\providecommand{\tightlist}{\setlength{\itemsep}{0pt}\setlength{\parskip}{0pt}}
\providecommand{\pandocbounded}[1]{#1}
\pagestyle{fancy}
\fancyhf{}
\fancyhead[L]{\sffamily\footnotesize\color{muted}OPEN PROBLEMS IN NUMERICAL LINEAR ALGEBRA}
\fancyhead[R]{\sffamily\bfseries\footnotesize\color{accent}MI-24}
\fancyfoot[L]{\parbox[b]{0.9\textwidth}{\sffamily\fontsize{8}{10}\selectfont\color{muted}Editorial ratings; a bounded search cannot certify that a problem remains unsolved.\\Verification check: 2026-09-18}}
\fancyfoot[R]{\sffamily\footnotesize\color{muted}\thepage}
\renewcommand{\headrulewidth}{0.3pt}
\renewcommand{\headrule}{\hbox to\headwidth{\color{accent}\leaders\hrule height \headrulewidth\hfill}}
\makeatletter
\renewcommand{\section}{\Needspace{5\baselineskip}\@startsection{section}{1}{0pt}{12pt plus 2pt minus 2pt}{4pt}{\sffamily\bfseries\large\color{ink}}}
\renewcommand{\subsection}{\Needspace{4\baselineskip}\@startsection{subsection}{2}{0pt}{9pt plus 1pt minus 1pt}{3pt}{\sffamily\bfseries\normalsize\color{ink}}}
\makeatother
\setcounter{secnumdepth}{0}
\begin{document}
{\sffamily\small\color{muted}Matrix inequalities and norms\par}
\vspace{3pt}
{\raggedright\hyphenpenalty=10000\exhyphenpenalty=10000\sffamily\bfseries\fontsize{21}{24}\selectfont\color{ink}Schatten
norm complement involving Lin's positive definite quantity\par}
\vspace{7pt}
{\sffamily\small\textbf{Difficulty:} challenging\quad\textbf{Importance:} interesting
to specialist\par}
{\sffamily\small\bfseries\color{accent}Status: Lean verified\par}
\vspace{4pt}
{\color{accent}\hrule height 0.8pt}
\vspace{6pt}
\textbf{Rating rationale (historical):} Extending the proved trace and
Frobenius cases to all Schatten norms is challenging; the comparison of
these particular means has specialist importance.

\subsection{Resolution --- 2026-09-11}\label{resolution-2026-09-11}

\textbf{Solved (affirmative).} George Stepaniants's
\href{https://github.com/ajt60gaibb/OpenProblemsInNLA/blob/main/matrix-inequalities-and-norms/MI-24/solution.md}{complete
proof}, \textbf{Theorem 1}, proves the displayed inequality for every
dimension, every complex positive definite pair and every
\(1\le p\le\infty\). The published Heron norm inequality and a positive
matrix comparison give \(\|A+B+2G\|_p\le\|A+B+G+L\|_p\). Since
\(2(A+B+G+L)=(A+B+2G)+(A+B+2L)\), the triangle inequality then proves
the desired comparison.
\href{https://github.com/ajt60gaibb/OpenProblemsInNLA/blob/main/matrix-inequalities-and-norms/MI-24/solution.pdf}{Proof
PDF} ·
\href{https://github.com/ajt60gaibb/OpenProblemsInNLA/blob/main/matrix-inequalities-and-norms/MI-24/solution.tex}{Standalone
TeX}.

The complete argument and its precise published input passed a
\href{https://github.com/ajt60gaibb/OpenProblemsInNLA/blob/main/references/stepaniants-mi24-2026-09-11/verification/reviews/MI-24-review.md}{separate
Codex-agent review}. It was developed with ChatGPT/Codex; that September
11 verification was independent agent review. The later Lean
verification is recorded below; no external human peer review is
claimed.
\href{https://github.com/ajt60gaibb/OpenProblemsInNLA/blob/main/references/stepaniants-mi24-2026-09-11/README.md}{Authorship,
source checks, reproduction and public-branch audit}. The contribution
is the convexity deduction from existing comparisons. The original ID,
path, statement, historical ratings and earlier status check below
remain intact.

\subsection{Lean proof and verification
evidence}\label{lean-proof-and-verification-evidence}

\textbf{The complete original target is Lean verified, 2026-09-18
(UTC).} The
\href{https://github.com/sgstepaniants/OpenProblemsInNLA/blob/97ff57c89775783bb7ebe952315840be86733417/matrix-inequalities-and-norms/MI-24/lean/Solution.lean}{immutable
proof} at revision \textit{97ff57c89775783bb7ebe952315840be86733417}
passed
\href{https://github.com/sgstepaniants/OpenProblemsInNLA/actions/runs/35310937025/job/105492723099}{non-root
Linux run 35310937025}. The
\href{https://github.com/ajt60gaibb/OpenProblemsInNLA/blob/main/matrix-inequalities-and-norms/MI-24/lean/verification/linux-2026-09-18/README.md}{retained
execution evidence} binds all 212 submitted project inputs and 22
exported statements. Real Comparator, default-kernel replay, standard
transitive axioms, sandbox checks and rejection controls passed.
Subsequent publication and PR merge checkouts are checked separately.

The main declaration \textit{NLA.MI24.schattenComplementConjecture}
proves the unchanged inequality for every complex positive definite
pair, every positive dimension, every real finite Schatten exponent at
least one, and the infinity endpoint. The original ordered formulas for
both matrix means are retained; no symmetry of Lin's quantity or
commutativity is assumed. The finite and infinity semantics theorems
establish the actual Schatten definitions. The needed Heron and Furuta
comparisons are proved internally. The
\href{https://github.com/ajt60gaibb/OpenProblemsInNLA/blob/main/matrix-inequalities-and-norms/MI-24/lean/Challenge.lean}{22
exact contracts} and
\href{https://github.com/ajt60gaibb/OpenProblemsInNLA/blob/main/matrix-inequalities-and-norms/MI-24/lean/IMPLEMENTATION-MAP.json}{implementation
map} identify the definitions, assumptions and proofs.

Convexity deduction and formalization: \textbf{George Stepaniants},
Department of Computing and Mathematical Sciences, California Institute
of Technology. The original conjecture remains attributed to Ghabries,
Abbas, Mourad and Assi; the Heron comparison to Dinh, Dumitru and
Franco; and the Furuta inequality and proof route to Furuta and Fujii.
Two independent statement reviews preceded implementation. Two
independent nonauthor final reviews approved the complete 34-module
source closure.
\href{https://github.com/ajt60gaibb/OpenProblemsInNLA/blob/main/matrix-inequalities-and-norms/MI-24/lean/reviews/README.md}{Review
records} disclose substantial OpenAI Codex assistance and distinguish
source review from executed verification.

The project pins Lean 4.33.1,
\href{https://github.com/leanprover-community/mathlib4/tree/0df444a360eaa60ab8c11dca51a86af692955474}{Mathlib}
and
\href{https://github.com/alerad/leancert/tree/621a43d7cf21f87872392a01e874f2f1dbddc926}{LeanCert}.
The proof consumes two kernel-mode LeanCert scalar bounds on one
interval of width one half; no matrix sampling or interval subdivision
is used. Every export has only \textit{propext},
\textit{Classical.choice} and \textit{Quot.sound} as transitive axioms.
Run \textit{lake build Solution} in
\textit{matrix-inequalities-and-norms/MI-24/lean};
\href{https://github.com/ajt60gaibb/OpenProblemsInNLA/blob/main/matrix-inequalities-and-norms/MI-24/lean/README.md}{reproduction
instructions} separately explain the actual local checks and the Linux
checker.

\subsection{Problem statement}\label{problem-statement}

For positive definite \(A,B\in\mathbb C^{n\times n}\), set

\[G=A^{1/2}(A^{-1/2}BA^{-1/2})^{1/2}A^{1/2},\qquad
L=A^{1/2}(B^{1/2}A^{-1}B^{1/2})^{1/2}A^{1/2}.\]

Here \(G=A\#B\), while \(L\) is the explicitly defined quantity used by
Lin; no symmetry of \(L(A,B)\) is assumed. Determine whether, for all
\(n\ge1\), all positive definite \(A,B\), and all \(1\le p\le\infty\),

\[\|A+B+G+L\|_p\le\|A+B+2L\|_p.\]

For \(p<\infty\),
\(\|X\|_p=(\mathop{\mathrm{tr}}\nolimits|X|^p)^{1/p}\), where
\(|X|=(X^*X)^{1/2}\); for \(p=\infty\), use the largest singular value.

Matrix means are central to interpolation of positive definite data. The
conjecture supplies a norm comparison between two concrete sums of
matrix functions; eigenvalue comparison between \(G\) and \(L\) alone is
not the statement being requested.

\subsection{References}\label{references}

\begin{enumerate}
\def\labelenumi{\arabic{enumi}.}
\tightlist
\item
  M. M. Ghabries, H. Abbas, B. Mourad and A. Assi, \emph{New
  log-majorization results concerning eigenvalues and singular values
  and a complement of a norm inequality}, preprint (2021), definition in
  §4, p.~11; Conjecture 4.1, p.~13.
  \href{https://arxiv.org/pdf/2105.13356}{PDF};
  \href{https://doi.org/10.1080/03081087.2022.2059050}{journal article}.
\item
  M. M. Ghabries, \emph{Contributions to Matrix Inequalities and Some
  Applications}, PhD thesis (2022), final Open Problems, Problem 5,
  pp.~109--110.
  \href{https://www.researchgate.net/publication/361793582_Contributions_to_Matrix_Inequalities_and_Some_Applications}{Author-uploaded
  thesis}.
\end{enumerate}

Status check (2026-09-10): the first source proves \(p=1,2\) and then
conjectures the full interval. The thesis explicitly retains that gap.
Current arXiv version v1, published metadata, the author's later July
2026 matrix-mean paper, and searches for ``Ghabries Conjecture 4.1
norm'', ``Lin complement geometric mean Schatten'', including 2025/2026,
did not reveal a proof or counterexample for the full interval. No
assertion of exhaustive coverage is made.
\end{document}
